\documentclass[11pt, a4paper]{article}

\usepackage[utf8]{inputenc}
\usepackage[T1]{fontenc}
\usepackage{amsmath, amssymb, amsthm}
\usepackage{graphicx}
\usepackage{hyperref}
\usepackage{cleveref}
\usepackage{enumitem}
\usepackage{booktabs}
\usepackage{geometry}
\usepackage{algorithm}
\usepackage{algpseudocode}

\geometry{margin=1in}

\title{\textbf{Sapien: A Decentralized Protocol for Verifiable, Consensus-Based Quality Signals in AI Workflows}}

\author{Chad Lynch, Rowan Stone, Trevor Koverko, Kelly Ryan \\
\href{mailto:chad@sapien.io}{chad@sapien.io} \\
Sapien Protocol}

\date{\today}

\begin{document}

\maketitle

\begin{abstract}
As Artificial Intelligence (AI) systems, particularly Large Language Models (LLMs), become increasingly ubiquitous, the integrity and quality of their training data and output validation become critical. Current methods for data labeling and reinforcement learning from human feedback (RLHF) rely heavily on centralized, opaque black-box services, lacking transparency and verifiable quality assurance. This paper presents Sapien, a decentralized "Proof-of-Quality" (PoQ) protocol implemented on the Ethereum Virtual Machine (EVM). Sapien introduces a mechanism for stake-weighted human verification of AI-generated content, utilizing a novel consensus engine that integrates financial staking with non-transferable reputation. We define the protocol's participant roles, formalize the commit-reveal consensus mechanism, and analyze the economic security properties of the system, including its resistance to Sybil attacks and stake centralization through quadratic weighting functions.
\end{abstract}

\section{Introduction}

The efficacy of modern AI models is fundamentally constrained by the quality of data used for training and fine-tuning. While model architectures have become increasingly standardized, the "data layer" remains fragmented and largely unverifiable. The current paradigm relies on centralized vendors for data annotation and Reinforcement Learning from Human Feedback (RLHF), creating a "trust me" environment where the provenance and quality of human labels are obfuscated.

This lack of transparency introduces significant risks:
\begin{enumerate}
    \item \textbf{Quality Degradation:} Without verifiable audits, labeling errors propagate into model behaviors.
    \item \textbf{Centralization Risks:} Dependence on a few large providers creates bottlenecks and censorship vectors.
    \item \textbf{Misaligned Incentives:} Workers in low-wage jurisdictions maximize throughput rather than accuracy, leading to "lazy labeling."
\end{enumerate}

We propose Sapien, a decentralized protocol that acts as a "Quality Oracle" for AI systems. Sapien leverages blockchain technology to create an immutable audit trail of human judgment. By aligning economic incentives through a staking-slashing mechanism and utilizing a consensus algorithm derived from quadratic voting research, Sapien enables a trustless market for high-fidelity data validation.

\section{System Architecture}

The Sapien protocol is architected as a modular suite of smart contracts in Solidity. The system facilitates the interaction between four primary agents: Originators, Contributors, Validators, and Oracles.

\subsection{Participant Roles}

\begin{itemize}
    \item \textbf{Originators:} Entities requiring verified data or agent supervision. They initialize projects, define quality criteria (Task Definition Specs), and fund reward pools.
    \item \textbf{Contributors:} Workers who perform the primary task (e.g., image annotation, prompt completion). They must stake tokens to claim work slots, providing "skin in the game."
    \item \textbf{Validators:} Independent reviewers who audit the work of Contributors. They reach consensus on quality scores.
    \item \textbf{Oracles (Adapters):} Middleware that bridges off-chain data (e.g., images stored in IPFS, CVAT interfaces) with on-chain coordination.
\end{itemize}

\subsection{Core Components}

The protocol logic is distributed across five primary contracts:
\begin{itemize}
    \item \texttt{SapienCore}: The central coordinator managing project lifecycles, contribution tracking, and finalization triggers.
    \item \texttt{ValidationOracle}: A stateless consensus engine that manages the commit-reveal scheme and validator queueing.
    \item \texttt{SapienTrust}: A non-transferable reputation system implementing the "Proof-of-Quality" (PoQ) score, which evolves based on historical performance.
    \item \texttt{SapienVault}: An ERC-4626 compliant staking vault that manages deposits, locks, and slashing.
    \item \texttt{Rewards}: Handles the escrow and distribution of ERC-20 reward tokens.
\end{itemize}

\section{The Proof-of-Quality Consensus Mechanism}

The core innovation of Sapien is its consensus mechanism, which validates the \textit{quality} of work rather than the ordering of transactions.

\subsection{Workflow Lifecycle}

The validation process follows a strict state machine:
\begin{enumerate}
    \item \textbf{Project Initialization:} Originator $O$ creates a project $P$ with a reward pool $R$ and a required validator count $k$.
    \item \textbf{Contribution:} Contributor $C$ stakes collateral $S_C$ and submits work $W$.
    \item \textbf{Commit Phase:} A set of validators $V = \{v_1, ..., v_k\}$ claim tasks. To prevent herding (where validators copy early responses), they execute a commit-reveal scheme. Each validator $v_i$ generates a score $s_i \in [0, 100]$ and a random salt $\rho_i$, submitting the hash:
    \begin{equation}
        H_i = \text{Keccak256}(s_i || S_{v_i} || \rho_i)
    \end{equation}
    where $S_{v_i}$ is the validator's staked capacity.
    \item \textbf{Reveal Phase:} After the commit window closes, validators reveal $(s_i, \rho_i)$. The contract verifies $H_i$.
    \item \textbf{Consensus Calculation:} The protocol calculates a weighted average score $\bar{S}$ and identifies outliers.
\end{enumerate}

\subsection{Weighting Algorithms}

Sapien utilizes pluggable consensus modules. The default and most robust algorithm, \textbf{Hybrid Consensus}, combats plutocracy (whale dominance) and Sybil attacks by weighing votes based on both stake and reputation.

The weight $w_i$ of validator $v_i$ is calculated as:
\begin{equation}
    w_i = \min\left( \sqrt{S_{v_i}} \cdot \text{Rep}_{v_i}, \quad 0.30 \cdot \sum_{j=1}^k w_j \right)
\end{equation}

Where:
\begin{itemize}
    \item $S_{v_i}$ is the staked amount. The square root function $\sqrt{S_{v_i}}$ introduces \textit{quadratic voting} properties, diminishing the marginal influence of additional capital.
    \item $\text{Rep}_{v_i}$ is the validator's PoQ reputation score.
    \item The $30\%$ cap ensures no single entity can dictate the consensus outcome, requiring a minimum of 4 distinct entities for valid consensus.
\end{itemize}

Other available algorithms include \textit{SqrtStake} ($w_i = \sqrt{S_{v_i}}$) and \textit{CappedLinear} ($w_i = \min(S_{v_i}, \text{Cap})$).

\subsection{Outlier Detection and Slashing}

To enforce honest behavior, the protocol employs statistical outlier detection. We define the consensus score $\mu$ and standard deviation $\sigma$ of the weighted votes.

A validator $v_i$ is classified as an outlier if:
\begin{equation}
    |s_i - \mu| > \max(15, 2\sigma)
\end{equation}

Outliers are subject to slashing—a punitive mechanism where a portion of their stake is seized. The slash magnitude $\lambda$ is a function of the deviation magnitude (measured in standard deviations, or Z-score):

\begin{equation}
    \lambda(Z) = 
    \begin{cases} 
      10\% & \text{if } 1.5 \le Z < 2 \\
      25\% & \text{if } 2 \le Z < 3 \\
      50\% & \text{if } 3 \le Z < 4 \\
      75\% & \text{if } 4 \le Z < 5 \\
      100\% & \text{if } Z \ge 5 
   \end{cases}
\end{equation}

Slashed funds are not burned but are effectively redistributed to the remaining honest stakers in the \texttt{SapienVault} by burning the offender's shares without removing the underlying assets, thereby increasing the exchange rate of the remaining shares.

\section{Security Analysis}

\subsection{Sybil Resistance}
While standard quadratic voting is susceptible to Sybil attacks (splitting stake into many small accounts), Sapien mitigates this via the \texttt{SapienTrust} reputation system. High-weight influence requires not just capital, but a history of accurate validation. Building reputation is a time-bound process that cannot be instantly purchased, making Sybil attacks economically inefficient.

\subsection{Whale Resistance}
The sublinear $\sqrt{S}$ weighting curve ensures that influence does not scale linearly with capital. For example, to achieve $10\times$ the voting weight of a base validator, an attacker would need $100\times$ the capital. Combined with the hard $30\%$ cap per validator, the system ensures broad committee participation.

\subsection{Commit-Reveal Integrity}
The two-phase commit-reveal mechanism is critical for preventing "lazy validation" or "herding," where validators wait for others to vote and simply copy the average to avoid slashing. By cryptographically binding the score to a hidden salt, the protocol ensures independent judgment. Validators who commit but fail to reveal are penalized by slashing their locked capacity, preventing griefing attacks on the protocol's liveness.

\section{Conclusion}

Sapien introduces a novel primitive for the AI infrastructure stack: a decentralized quality layer. By combining cryptoeconomic incentives with a rigorous statistical consensus engine, Sapien provides a trustless mechanism for verifying AI data and behaviors. This protocol enables a shift from "trust-based" centralized labeling to "verification-based" decentralized coordination, essential for the safe and reliable deployment of autonomous AI agents.

\end{document}
